% ===================================================================
%  TABELL Nr 12 — Stationen. — Järnvägarna. — Fartygen.
%  Traduction 2026 d'apres texto/io, controlee sur texto/fr et
%  texto/es. Consigne : voir texto/sv/10-tabell-01.tex.
% ===================================================================

%%K t12-apar-1 apar
\VUsaut{4.0mm}
\VUtitre{15.0pt}{60}{TABELL Nr 12}
\VUsaut{2.4mm}
\VUpk{11.4pt}{40}{Stationen. --- Järnvägarna. --- Fartygen.}
\VUsaut{3.4mm}
\textsc{Anm.} --- \textit{De tidtabeller som används i varje land är olika, vi använder alltså
som exempel timmarna i den} \VUgras{franska tabellen.}

%%K t12-01-1 p
1. --- Jag har just rest genom Tyskland. Före avresan köpte jag en
\VUgras{tidtabellsbok} \textsuperscript{(15)} för att veta tågens avgångstid.
Sedan jag hade konstaterat att det lämpligaste tåget skulle avgå från Paris klockan nio på
kvällen, och komma till Strasbourg klockan ett och tretton minuter, förberedde jag min resa,
och dagen därpå lät jag föra mig till stationen.

%%K t12-02-1 p
2. --- När jag gick in i den stora avgångshallen, bakom en gammal
\VUgras{lantpräst} \textsuperscript{(2)}, som bar sin \VUgras{resväska}
\textsuperscript{(3)} i handen, hade redan många \VUgras{resande} \textsuperscript{(1)}
kommit.

%%K t12-02-2 p
Eftersom jag ville skicka ett \VUgras{telegram} \textsuperscript{(5)}, lät jag en
\VUgras{kontrollör} \textsuperscript{(6)} visa mig \VUgras{telegrafkontoret}
\textsuperscript{(4)}.

%%K t12-03-1 p
3. --- Innan jag gick in i \VUgras{väntsalen} \textsuperscript{(7)} gick jag och försåg mig
med en \VUgras{biljett} \textsuperscript{(9)} tur och retur.

%%K t12-04-1 p
4. --- Jag väntade inte länge vid \VUgras{luckan} \textsuperscript{(8)}, för det var få
personer där. En \VUgras{soldat} \textsuperscript{(10)}, som reste på permission, var så artig
att lämna mig sin plats i kön, så att jag hade tillräcklig tid att be om några upplysningar av
\VUgras{stationsinspektorn} \textsuperscript{(13)}, som samtalade med \VUgras{kommissarien}
\textsuperscript{(12)} för uppsikten och med den tjänstgörande
\VUgras{gendarmen} \textsuperscript{(11)}.

%%K t12-tit-1 sub
\VUpk{10.2pt}{40}{Inskrivningen av bagaget.}
\VUsaut{3.0mm}

%%K t12-05-1 p
5. --- Sedan jag hade köpt en tidning i \VUgras{bokhandeln} \textsuperscript{(14)}, lät jag
väga mitt \VUgras{bagage} \textsuperscript{(17)} som en \VUgras{bärare} \textsuperscript{(16)}
just hade fört dit på en
\VUgras{packvagn} \textsuperscript{(19)}; \VUgras{tjänstemannen}
\textsuperscript{(22)} som ansvarar för \VUgras{vågen} \textsuperscript{(21)} meddelade mig,
att min koffert väger trettiofem kilogram; det gjorde en
\VUgras{övervikt} på fem kilogram; för vilken jag är skyldig en tilläggsavgift,
vid bagagets \VUgras{lucka} \textsuperscript{(23)}.

%%K t12-06-1 p
6. --- Eftersom jag var för tidig, hade jag ledig tid att betrakta de resandes rörelse på
stationen. Somliga lämnade in eller hämtade ut sitt handbagage vid
\VUgras{(bagage)inlämningen} \textsuperscript{(25)}; andra rådfrågade
\VUgras{tidtabellerna} \textsuperscript{(26)}. De ankommande resande skyndade
till \VUgras{utgången} \textsuperscript{(32)} med sin \VUgras{kappsäck} \textsuperscript{(24)}
i handen. Några väntade vid \VUgras{(bagage)utlämningen} \textsuperscript{(27)} och visade sin
\VUgras{polett} \textsuperscript{(29)} för att få tillbaka sina koffertar, \VUgras{lådor}
\textsuperscript{(28)} eller
\VUgras{korgar} \textsuperscript{(30)}.

%%K t12-07-1 p
7. --- \VUgras{Accistjänstemännen} \textsuperscript{(33)} undersökte omsorgsfullt packen för
att fastställa om man har något att förtulla.

%%K t12-08-1 p
8. --- Somliga resande gick till \VUgras{restaurationen} \textsuperscript{(34)}, där en
\VUgras{kypare} \textsuperscript{(35)} bemödade sig att betjäna dem. De måste skynda sig, för
\VUgras{tåget} \textsuperscript{(36)}, som är ett snälltåg, ska inte stanna länge på
stationen.

%%K t12-09-1 p
9. --- Sedan gick jag till avgångsperrongen och sökte en genomgående
\VUgras{vagn} \textsuperscript{(37)} för att inte vara tvungen att byta vagn
under resan. Jag steg upp i en korridorvagn som innehåller en \VUgras{kupé}
\textsuperscript{(38)} för rökare och en kupé reserverad för « ensamma damer ».

%%K t12-tit-2 sub
\VUpk{10.2pt}{40}{Avresa under natten.}
\VUsaut{3.0mm}

%%K t12-10-1 p
10. --- Sedan jag hade stigit upp på \VUgras{fotsteget} \textsuperscript{(40)}, stängt
\VUgras{dörren} \textsuperscript{(39)}, rullat upp \VUgras{gardinen} \textsuperscript{(41)}
och lagt mitt handbagage på \VUgras{nätet} \textsuperscript{(43)}, satte jag mig på
\VUgras{bänken} \textsuperscript{(42)}. När jag såg \VUgras{lampskötaren}
\textsuperscript{(44)} tända lamporna, tänkte jag, att jag skulle behöva tillbringa en natt i
vagnen och jag kallade på den tjänsteman som ansvarar för uthyrningen av \VUgras{huvudkuddar}
\textsuperscript{(45)} för att be om en av dem.

%%K t12-11-1 p
11. --- Tågets konduktör kom och kontrollerade biljetterna. Jag frågade honom varför vi ännu
inte hade avgått, eftersom det redan var dags. Han svarade mig att
\VUgras{tågchefen} \textsuperscript{(46)} håller på att låta sätta in cyklar i
\VUgras{resgodsvagnen} \textsuperscript{(47)}. Dessutom hade man ännu inte bytt
alla fotvärmare \textsuperscript{(48)}.

%%K t12-12-1 p
12. --- Men vi skulle snart avgå, för \VUgras{eldaren} \textsuperscript{(50)}, på sin
\VUgras{tender} \textsuperscript{(49)} tog kol för att egga
\VUgras{lokomotivets} \textsuperscript{(54)} eld, och \VUgras{maskinisten}
\textsuperscript{(51)} väntade bara på tecknet från \VUgras{underinspektorn}
\textsuperscript{(52)}, som stod på \VUgras{perrongen} \textsuperscript{(53)}.

%%K t12-13-1 p
13. --- Lokomotivets \VUgras{panna} \textsuperscript{(56)} mullrade dovt;
\VUgras{skorstenen} \textsuperscript{(55)} utspydde strömmar av rök blandad med
\VUgras{ånga} \textsuperscript{(60)}. En gäll \VUgras{vissling}
\textsuperscript{(59)} ljöd och \VUgras{kolvarna} \textsuperscript{(58)} började röra sig, och
drev den tunga maskinens \VUgras{hjul} \textsuperscript{(57)}.

%%K t12-tit-3 sub
\VUsaut{3.0mm}
\VUpk{10.2pt}{40}{Avfärd!}
\VUsaut{3.0mm}

%%K t12-14-1 p
14. --- Tågets fart, medan det gick på \VUgras{rälsen} \textsuperscript{(64)}, ökades;
\VUgras{perrongerna} \textsuperscript{(65)} försvann; vi passerade
\VUgras{avträdena} \textsuperscript{(66)}, och också en rad vagnar som var
uppställda på det andra \VUgras{spåret} \textsuperscript{(63)}:
\VUgras{sovvagnar} \textsuperscript{(67)}, en \VUgras{restaurangvagn}
\textsuperscript{(68)}, en \VUgras{postvagn} \textsuperscript{(69)}.

%%K t12-noto-1 noto
\VUnotes{143.2mm}{%
\textit{(*) Anm.} --- \textit{Självfallet är beskrivningen av resan enligt gränsen före
kriget.}%
}

%%K t12-15-1 p
15. --- Sedan gled \VUgras{godsstationen} \textsuperscript{(71)} förbi våra ögon. Under
skjulen lastade tjänstemän \VUgras{godset} \textsuperscript{(72)} som sänds med godståg.
\VUgras{Växlingskarlar} \textsuperscript{(76)} vid
\VUgras{vattenreservoaren} \textsuperscript{(73)} växlade \VUgras{kreatursvagnar}
\textsuperscript{(75)} och \VUgras{plattformsvagnar} \textsuperscript{(79)} för att bilda ett
\VUgras{godståg} \textsuperscript{(80)}.

%%K t12-16-1 p
16. --- Vi passerade den sista \VUgras{växeln} \textsuperscript{(78)}, vid vilken växelkarlen
stod; vi for förbi \VUgras{signalskivan} \textsuperscript{(86)} och stationen försvann i
fjärran.

%%K t12-17-1 p
17. --- Snart for vi över en \VUgras{järnbro} \textsuperscript{(74)} som går över
\VUgras{floden} \textsuperscript{(81)}. Sedan vi hade farit över den bron, följde
vi älven, på vilken kraftiga \VUgras{bogserbåtar} \textsuperscript{(82)} drog tunga
\VUgras{pråmar} \textsuperscript{(83)}. Vi såg dessutom ett
\VUgras{passagerarfartyg} \textsuperscript{(84)}, när det lade till vid en
\VUgras{ponton} \textsuperscript{(85)}.

%%K t12-18-1 p
18. --- Sedan vi rasande snabbt hade passerat flera stationer, for vi genom några
\VUgras{tunnlar} \textsuperscript{(87)} och lämnade bakom oss otaliga
\VUgras{små hus} \textsuperscript{(88)} åt \VUgras{banvakter}. Vaggad av tågets
skakningar somnade jag djupt.

%%K t12-tit-4 sub
\VUsaut{3.0mm}
\VUpk{10.2pt}{40}{Stationen Deutsch-Avricourt.}
\VUsaut{3.0mm}

%%K t12-19-1 p
19. --- Jag väcktes plötsligt av en tjänstemans röst som ropade: « Deutsch-Avricourt!
\textsuperscript{(*)}. Alla stiger av! » Tullens undersökning och alla gränsens tröttsamma
formaliteter ägde rum. Jag måste ställa mitt fickur efter stationens \VUgras{stora ur}
\textsuperscript{(89)}, som gick femtiofem minuter för fort; knappt vaken, urskilde jag med
svårighet den
\VUgras{stora visaren} \textsuperscript{(90)} från den \VUgras{lilla}
\textsuperscript{(91)}; själva \VUgras{urtavlan} \textsuperscript{(92)} var alldeles förvirrad
av morgondimman.

%%K t12-20-1 p
20. --- Medan tulltjänstemännen gjorde sin undersökning, tydde jag gäspande de
\VUgras{affischer} \textsuperscript{(93)} som pryder stationens murar. Jag åt
frukost för att fördriva tiden, rådfrågade kartan över det tyska \VUgras{nätet}
\textsuperscript{(94)} för att se hur långt avstånd som ännu skiljer mig från Strasbourg, och
jag gick in i min vagn igen, stärkt av hoppet att snart nå målet för min resa.
\VUsaut{6.0mm}
\VUfilet{22mm}
